\section{Model, prediction convention, and stability restriction}
\label{sec:theory-model}

Let $\mathcal X=\{-1,1\}^d$, let $D$ be any distribution on $\mathcal X$,
and let $f:\mathcal X\to\{-1,1\}$ be the target. A network has $H=L-1$
hidden layers, each of width $N$, tanh activations, and a linear readout:
\[
 h^1=\tanh(W^1\bar x),\qquad
 h^s=\tanh(W^s h^{s-1})\ (2\le s\le H),\qquad
 z_w(x)=v^\top h^H,
 \qquad \bar x=(x,1).
\]
The first layer includes a bias coordinate; internal constants are carried
by ordinary hidden wires. Let $\mu_{N,t}$ be the product Gaussian with entry
variances $t^2/d$ in $W^1$, $t^2/N$ in the internal matrices, and $t^2/N^2$
in $v$, where $0<t\le1$. The exponent convention
$t=N^{(1-\gamma)/2}$ therefore gives every $\gamma\ge1$. Define
\[
 \kappa_N(w)=\frac dN\|W^1\|_F^2+
           \sum_{s=2}^H\|W^s\|_F^2+N\|v\|^2,
 \qquad E_v(w)=N\|v\|^2.
\]
In these energy coordinates the Gaussian density is proportional to
$\exp[-N\kappa_N/(2t^2)]$.

A disturbance is added after every hidden activation and at the readout:
\[
 h^1_\eta=\tanh(W^1\bar x)+\eta^1,\qquad
 h^s_\eta=\tanh(W^s h^{s-1}_\eta)+\eta^s,\qquad
 z_w(x;\eta)=v^\top h^H_\eta+\eta^{\rm out}.
\]
The simultaneous disturbance budget is
\begin{equation}\label{eq:standalone-noise}
 \|\eta\|_*^2=\sum_{s=1}^H\frac{\|\eta^s\|_2^2}{N}
                       +|\eta^{\rm out}|^2.
\end{equation}
This is an adversarial finite perturbation, not a Gaussian-noise average or
a derivative evaluated only at the clean computation. Noisy hidden states
need not stay in $[-1,1]^N$.

Our training likelihood is
\begin{equation}\label{eq:standalone-likelihood}
 \mathcal L_S(w)=\prod_{i=1}^n
       \mathbf1\{f(X_i)z_w(X_i)>1/2\}.
\end{equation}
At test time we count $f(x)z_w(x)\le1/4$ as an error. This stronger loss
also bounds sign error. There is no requirement that unobserved inputs have
any clean output margin.

Clip the output at the training margin:
\[
 T(u)=\max\{-1/2,\min\{1/2,u\}\}.
\]
For a calibration parameter $c>0$, set
\begin{equation}\label{eq:standalone-radius}
 r_c(w)=\frac{c}{\sqrt{1+H E_v(w)}}.
\end{equation}
Let $\mathcal R_{N,c}$ be the set of networks satisfying
\begin{equation}\label{eq:standalone-support}
 |T(z_w(x;\eta))-T(z_w(x))|\le1/16
 \quad\text{for every }x\in\mathcal X
 \text{ and every }\|\eta\|_*\le r_c(w).
\end{equation}
The prior is exactly $\mu_{N,t}$ conditioned on $\mathcal R_{N,c}$.
The posterior $Q_{S,N,t,c}$ additionally conditions on
\eqref{eq:standalone-likelihood}. Both conditioning events have positive
probability by the target construction below. No tractable sampling
algorithm is assumed or supplied by the theorem.

The clipping scale is set by the existing likelihood convention. For
example, changing a raw output from $10$ to $1$ causes no score change.
Clipping at $\pm1/2$ gives a weaker condition than clipping at $\pm1$ with
the same tolerance and noise radius, because
$T=T\circ T_1$, where $T_1$ clips at $\pm1$, and $T$ is $1$-Lipschitz.

\section{The learning theorem}
\label{sec:theory-main}

A width-$C$, depth-$H$ layered Boolean circuit has at most $C$ Boolean wires
at each of its $H$ stages. We use NAND gates, copies, and negations;
constant wires required by internal gate implementations are carried
through the circuit and count toward $C$. This is a universal Boolean
basis. Compiling another fixed basis into this one may change the width
and depth, and those resources must be counted. Its output is one final wire.
The circuit's logical depth must fit the $H$ hidden layers; idle stages
can be filled with copies.

\begin{theorem}[Polynomial depth at every width and every $\gamma\ge1$]
\label{thm:standalone-clipped}
There is a universal polynomial $P$ with the following property. Let $f$
be computed by such a width-$C$, depth-$H$ Boolean circuit, and let $D$ be
any input distribution. For every $c>0$ and every $0<\varepsilon,\delta<1$, if
\[
 n\ge P\bigl(d,L,C,c+c^{-1},\varepsilon^{-1},\log(1/\delta)\bigr),
\]
then with probability at least $1-\delta$ over the iid labeled sample,
\begin{equation}\label{eq:standalone-risk}
 \mathbb E_{w\sim Q_{S,N,t,c},\,X\sim D}
       \mathbf1\{f(X)z_w(X)\le1/4\}\le\varepsilon
\end{equation}
simultaneously for every $N\ge C$ and $0<t\le1$.
The polynomial degree is independent of width, depth, and Gaussian scale.
In particular, the result includes the actual finite-variance mean-field
scale $\gamma=1$ and every stronger-prior scale $\gamma>1$.
\end{theorem}

The same theorem holds if the prior is additionally conditioned on
$|z_w(x)|>1/2$ at every input, as in the original globally margin-conditioned
model. This variant is argued here; it is not part of the Lean-checked
statement, which concerns the robust-only posterior. The target cylinder constructed below already satisfies this extra
condition. Its probability bound and every subsequent statistical step are
therefore unchanged. Under that additional condition, test-margin error
coincides with sign error. Dropping the extra condition is permitted, but is
not required for the result.

The width threshold concerns a layered Boolean circuit, not an arbitrary
real-valued neural-network representation of the same Boolean function.
The theorem does not assert that every function with a small real neural
circuit has a comparably small Boolean circuit at the same depth, or that
all its low-energy neural representations satisfy the stability condition.
The proof instead supplies one robust implementation of each Boolean
circuit, together with enough Gaussian probability around that implementation.

Define the posterior correct-margin probability
$p_S(x)=Q_{S,N,t,c}\{f(x)z_w(x)>1/4\}$.
Equation~\eqref{eq:standalone-risk} immediately implies
\[
 D\{x:p_S(x)\le2/3\}\le3\varepsilon,
 \qquad D\{x:p_S(x)\le1/2\}\le2\varepsilon.
\]
The same bounds hold for posterior sign correctness, since the correct-sign
event contains the correct-margin event. These conclusions allow arbitrary
posterior mixtures; they do not require concentration on one predictor.

\section{Compression and the statistical proof}
\label{sec:theory-statistical}

The constructive compression statement needed here concerns numerical
outputs before any robustness condition is used.

\begin{lemma}[Entropy compression before the prediction comparison]
\label{lem:raw-entropy-compression}
For $d\ge1$, $L\ge2$, $K\ge1$, $0<r\le1$, and
$0<\alpha\le1/8$, there is a finite class $\mathcal H$ of real-valued
bounded-weight depth-$L$ tanh networks, independent of $N,t$, such that
\[
 \log|\mathcal H|\le
       \operatorname{poly}(d,L,K,r^{-1},\alpha^{-1}).
\]
If $Q\ll\mu_{N,t}$ satisfies
$D_{\rm KL}(Q\Vert\mu_{N,t})\le NK/t^2$, there is one distribution
$\Gamma_Q$ on $\mathcal H$ with the following property. For every fixed
$x$, one can couple $w\sim Q$ and $h\sim\Gamma_Q$ so that, except with
probability at most $\alpha$, there is a disturbance $\eta$ with
\[
 \|\eta\|_*<r/2,\qquad
 |h(x)-z_w(x;\eta)|\le r/8.
\]
The couplings and disturbances may depend on $x$, but their two marginals
are the same at every input. No robustness or clean-margin assumption is
made on $Q$.
\end{lemma}

The proof is constructive: it retains a bounded number of input directions
carrying conditional Gaussian relative entropy in each layer, replaces the
remaining directions by Gaussian randomness, samples a bounded number of
source neurons, and rounds a bounded-size network to a fixed grid.
The resulting virtual computation is the original network under small
postactivation disturbances. The appendix gives the complete construction;
the prediction comparison below is the first step that uses stability.

\begin{lemma}[Readout truncation]
\label{lem:standalone-readout-tail}
Suppose a posterior is the raw Gaussian conditioned on an event of mass at
least $\exp(-ND_*/t^2)$, where $D_*\ge1$. For $0<\beta\le1/2$, put
\[
 B=\max\{8,8[D_*+\log(1/\beta)]\}.
\]
It assigns mass at most $\beta$ to $E_v>B$. Its conditional law $Q'$ on
$E_v\le B$ satisfies
$D_{\rm KL}(Q'\Vert\mu_{N,t})\le N(D_*+1)/t^2$.
\end{lemma}
\begin{proof}
Under the Gaussian, $NE_v/t^2$ has the chi-square distribution with $N$
degrees of freedom. Its exponential moment gives, for $B\ge8$,
\[
 \mu_{N,t}\{E_v>B\}
 \le e^{-NB/(4t^2)}2^{N/2}
 \le e^{-NB/(8t^2)}.
\]
Divide by the conditioning-event probability and use $N/t^2\ge1$.
The further conditioning event has raw probability at least
$(1-\beta)e^{-ND_*/t^2}$, so its negative logarithm is at most
$ND_*/t^2+\log[1/(1-\beta)]\le N(D_*+1)/t^2$.
\end{proof}

\begin{proof}[Statistical proof of Theorem~\ref{thm:standalone-clipped}]
The next section constructs a subset of the support with correct noisy
margin on every input and raw mass at least $e^{-ND_*/t^2}$, where
$D_*=\operatorname{poly}(d,L,C,c)$ is independent of $N,t$.
It is contained in the posterior conditioning event for every correctly
labeled sample. Apply Lemma~\ref{lem:standalone-readout-tail} with
$\beta=\varepsilon/4$, and write $Q'$ for the retained law. Set
\[
 K=D_*+1,\qquad
 r=\min\{1/6,c/\sqrt{1+HB}\},\qquad
 \alpha=\varepsilon/12.
\]
Every retained network satisfies the score-stability condition at the common
radius $r$. Lemma~\ref{lem:raw-entropy-compression} and the $1$-Lipschitz
property of $T$ therefore give pointwise couplings with
\[
 \Pr\{|T(h(x))-T(z_w(x))|>1/12\}\le\alpha,
\]
because $1/16+r/8\le1/12$.

Use the bounded ramp loss
\[
 \psi(u)=
 \begin{cases}
 1,&u\le1/3,\\
 5-12u,&1/3<u<5/12,\\
 0,&u\ge5/12.
 \end{cases}
\]
At a training input, $f(x)z_w(x)>1/2$ implies $f(x)T(z_w(x))=1/2$.
A successful coupling therefore gives $f(x)T(h(x))\ge5/12$, so its ramp
loss is zero. The candidate's averaged empirical ramp risk is at most
$\alpha$. At a test input with $f(x)z_w(x)\le1/4$, a successful coupling
gives $f(x)T(h(x))\le1/3$, so its ramp loss is one.

For any finite class of $[0,1]$-valued losses, a relative Chernoff bound
and a union bound give, with probability at least $1-\delta$,
\[
 \mathbb E_D\psi(f(X)T(h(X)))
 \le\frac2n\sum_{i=1}^n\psi(f(X_i)T(h(X_i)))
       +\frac{8[\log|\mathcal H|+\log(1/\delta)]}{n}
 \quad(h\in\mathcal H).
\]
For completeness, $e^{-u}\le1-(1-e^{-1})u$ on $[0,1]$ gives the usual
lower-tail exponent at least the mean divided by eight. This proves the
displayed bound for one loss, and a union bound proves simultaneity.
Average under the sample-dependent law $\Gamma_{Q'}$, apply the population
comparison, and restore the discarded posterior mass. The result is
\[
 \Pr_{Q_{S,N,t,c},D}\{f(X)z_w(X)\le1/4\}
 \le\beta+3\alpha+
       \frac{8[\log|\mathcal H|+\log(1/\delta)]}{n}.
\]
Choose $n$ so the last term is at most $\varepsilon/2$.
Every quantity entering $\log|\mathcal H|$ is bounded by a universal
polynomial in the theorem's parameters. Its class and the displayed sample
event are independent of $N,t$, establishing the simultaneous statement.
\end{proof}

\section{A robust target cylinder at every width}
\label{sec:theory-target-cylinder}

We give the probability bound used in the statistical proof. It uses a
robust digital implementation, not a claim that every implementation is
stable. Let $N\ge C$. Partition the neurons into $C$ balanced clone groups
of sizes $k_a\in\{\lfloor N/C\rfloor,\lceil N/C\rceil\}$, so
$C/2\le N/k_a\le2C$. Let $P$ project onto vectors constant on each clone
group. Fill unused logical slots with available constants or copies of an
existing Boolean wire. The ideal wire
vector $b^s(x)$ then has entries in $\{-1,1\}$, belongs to the range of
$P$, and has Euclidean norm $\sqrt N$.

Use gain $A$ in the tanh simulation of every Boolean gate. For example,
NAND on $u,v\in\{-1,1\}$ is the sign of $1-u-v$; copies and negations
use $u$ and $-u$. Bias constants in internal layers are supplied by counted
constant wires. The ideal signed preactivation of every gate is at least
$A$. In the clone implementation, each coefficient from a logical wire is
divided equally among that wire's clones. Set
\begin{align*}
 e_*&=\frac1{16\sqrt{2C}},&
 \tau_0&=\frac{e_*}{4},&
 b_L&=\sqrt{16\log9+8\log(4L)},\\
 A&=\max\{32b_L,2\log(8/e_*)\},&
 a&=\max\{2,256c\},&
 \rho&=\min\{1/4,A/[4(d+1)],Ae_*/(64\sqrt C)\}.
\end{align*}
The teacher readout averages the clones of the output wire with total
coefficient $a$.

In Gaussian coordinates $\sqrt d\,W^1$, $\sqrt N\,W^s$, and $Nv$, all
entries have variance $t^2$. Restrict each first-layer coordinate, each
internal row's projection onto every normalized clone-group direction, and
each readout coordinate to an interval of radius $\rho t$ about its teacher
value. The internal matrices then decompose as
\[
 W^s=W_*^s+R^s+E^s,\qquad
 R^s=R^sP,\quad \|R^s\|_{\rm op}\le\rho t\sqrt C,
 \quad E^sP=0.
\]
The remaining Gaussian bulk $E^s$ is independent of the constrained
coordinates. A $1/4$-net of the unit sphere has at most $9^N$ points,
and the usual two-net Gaussian bound gives
\[
 \Pr\{\|E^s\|_{\rm op}>b_Lt\}
 \le2\exp(2N\log9-Nb_L^2/8)\le\frac1{2L}.
\]
Thus all bulk norms are at most $b_Lt$ with conditional probability at
least $1/2$.

\paragraph{Finite disturbance robustness.}
Allow a separate RMS disturbance of size $\tau_0$ at each hidden layer and
scalar output disturbance at most $\tau_0$. This includes the entire
total-budget ball $\|\eta\|_*\le\tau_0$.
Write $e_s=\|h^s_\eta-b^s\|_2/\sqrt N$.
If $e_{s-1}\le e_*$, every clone-group mean differs from its ideal bit by
at most $e_*\sqrt{N/k_a}\le1/16$. At most three teacher coefficients,
including a constant coefficient, enter a gate, so its intended signed
preactivation remains at least $A/2$.
For $\epsilon\in\{-1,1\}$ and real $u$, the elementary global secant bound
\[
 \epsilon q\ge A/2
 \quad\Longrightarrow\quad
 |\tanh(q+u)-\epsilon|\le2e^{-A/2}+8|u|/A
\]
holds: for $|u|<A/4$ use saturation, and otherwise the linear term is at
least two. Since $E^sb^{s-1}=0$ and
$\|h^{s-1}_\eta\|_2/\sqrt N\le1+\tau_0$, it follows that
\[
 e_s\le2e^{-A/2}+
          (8/A)\rho\sqrt C(1+\tau_0)
          +(8b_L/A)e_{s-1}+\tau_0\le e_*.
\]
Each of the four terms is at most $e_*/4$.
The first layer initializes the same invariant: its coordinatewise
preactivation perturbation is at most $\rho(d+1)\le A/4$, and its
postactivation disturbance contributes at most $\tau_0$ in RMS norm.
Consequently the noisy signed readout is at least
\[
 f(x)z_w(x;\eta)
 \ge\frac{15a}{16}-\rho(1+\tau_0)-\tau_0>1
 \qquad(x\in\mathcal X,\ \|\eta\|_*\le\tau_0).
\]

The readout center satisfies $\sqrt N\|v_*\|\ge a\sqrt{C/2}$, and its
permitted perturbation has $\sqrt N\|v-v_*\|\le\rho$.
Therefore $E_v\ge a^2C/8$, whence
\[
 r_c(w)\le\frac{2\sqrt2c}{a\sqrt C}
       \le\frac1{64\sqrt{2C}}=\tau_0.
\]
All clean and noisy clipped scores on the calibrated ball equal $f(x)/2$.
The entire cylinder, with the bulk-norm event imposed, therefore belongs
to $\mathcal R_{N,c}$ and satisfies the training likelihood for every sample.
This also proves why every fixed $c>0$ is allowed: a larger calibration
requires only a polynomially larger readout coefficient.

\paragraph{Gaussian mass.}
There are $Nq$ constrained independent scalar coordinates, where
$q=d+2+(L-2)C$. The sum of their squared centers is at most $NB_0$, with
\[
 B_0=9dA^2+18(L-2)CA^2+2a^2C.
\]
These loose constants cover the fan-in-two gate implementation above.
For $Z\sim\mathcal N(0,t^2)$, $\rho\le1/4$, and any center $m$,
\[
 \Pr\{|Z-m|\le\rho t\}
 \ge\frac{2\rho}{e\sqrt{2\pi}}\exp(-m^2/t^2).
\]
Multiplication over the coordinates and the independent bulk event gives
\begin{equation}\label{eq:standalone-cylinder-mass}
 \mu_{N,t}(\text{target cylinder})
 \ge\exp(-ND_*/t^2),\qquad
 D_*=B_0+q\log\frac{e\sqrt{2\pi}}{2\rho}+\log2.
\end{equation}
Here $D_*\ge1$ and is bounded by a universal polynomial in $d,L,C,c$.
The construction and estimate hold at every integer width $N\ge C$.
